% ============================================================================
\section{紧化：$X(\Gamma)$ 与尖点}
\label{sec:compactification}
\label{sec:compactification-X}
% ============================================================================

\textbf{上一节留下了什么缺口？}

我们已经把 $Y(\Gamma)=\Gamma\backslash\HH$ 变成了黎曼面，
但它还不是紧的。
几何上说，你总可以沿着某个尖点方向越走越高；
解析上说，这意味着许多漂亮的紧曲面工具——尤其是 Riemann--Roch——还不能直接上场。
所以我们接下来的目标非常明确：把这些"漏掉的无穷远点"补回来。

\textbf{核心想法是把尖点当作真正的复点加入进去。}
在尖点附近，$\tau$ 本身不是好坐标，
因为平移 $\tau\mapsto\tau+h$ 会把不同的 $\tau$ 识别为同一点。
但指数坐标
\[
  q_h=e^{2\pi i\tau/h}
\]
会把这种平移变成同一个值；
于是尖点附近的条带就会变成一个穿孔圆盘。
把穿孔中心补上，尖点就成了普通的复点。

\textbf{这样做的收益是什么？}
我们得到紧黎曼面
\[
  X(\Gamma)=\Gamma\backslash\HH^*=\Gamma\backslash\bigl(\HH\cup\mathbb{P}^1(\Q)\bigr),
\]
其中新加进去的点就是尖点。
以后谈亏格、微分形式和模形式的几何解释时，真正起作用的对象都是这个紧化后的 $X(\Gamma)$。

\begin{definition}[尖点宽度与局部参数]
\label{def:cusp-width}
设 $s\in\mathbb{P}^1(\Q)$ 是 $\Gamma$ 的一个尖点，取
$\sigma\in\SL_2(\Z)$ 满足 $\sigma\infty=s$。
则
\[
  \sigma^{-1}\Gamma\sigma
\]
是 $\infty$ 的一个稳定子群，其中存在唯一正整数 $h$，使得
\[
  \pm\begin{pmatrix}1&h\\0&1\end{pmatrix}\in \sigma^{-1}\Gamma\sigma,
  \qquad
  \pm\begin{pmatrix}1&m\\0&1\end{pmatrix}\notin \sigma^{-1}\Gamma\sigma
  \quad (0<m<h).
\]
这个 $h$ 称为尖点 $s$ 的\textbf{宽度}（width）。
在该尖点附近的局部参数取为
\[
  q_s=q_h=e^{2\pi i\sigma^{-1}\tau/h}.
\]
当 $s=\infty$ 时，这就是 $q_h=e^{2\pi i\tau/h}$。
\end{definition}

\textbf{为什么这个坐标自然？}
若 $\tau\sim \tau+h$，则
\[
  e^{2\pi i(\tau+h)/h}=e^{2\pi i\tau/h}e^{2\pi i}=e^{2\pi i\tau/h}.
\]
也就是说，尖点方向上的平移对 $q_h$ 来说完全看不见。
因此 $q_h$ 恰好是商空间在尖点附近真正能"看见"的坐标。

% figures/ch02-cusp-q-coordinate.tex
% 尖点附近的条带通过 q_h=e^{2\pi i\tau/h} 映到小穿孔圆盘
\begin{figure}[ht]
\centering
\begin{tikzpicture}[>=Stealth, font=\small, scale=1.0]
  % strip
  \draw[thick] (-4.5,-1.7) rectangle (-1.5,1.7);
  \fill[blue!6] (-4.5,-1.7) rectangle (-1.5,1.7);
  \draw[dashed] (-4.5,0) -- (-1.5,0);
  \node[left] at (-4.55,0) {$T$};
  \node[above] at (-3.0,1.8) {$\{\,|\Re\tau|\le h/2,\ \Im\tau>T\,\}$};
  \draw[<->] (-4.3,-1.95) -- (-1.7,-1.95);
  \node[below] at (-3.0,-1.95) {宽度 $h$};
  \draw[->, thick, green!50!black] (-3.0,1.2) -- (-3.0,2.0);
  \node[green!50!black, right] at (-2.95,1.8) {$\Im\tau\to\infty$};
  \node[below] at (-3.0,-2.35) {在条带中，$\tau\sim \tau+h$};

  % arrow
  \draw[->, very thick] (-0.9,0) -- (0.7,0);
  \node[above, align=center] at (-0.1,0.1)
    {$q_h=e^{2\pi i\tau/h}$\\把平移变成同一个点};

  % punctured disk
  \draw[thick] (3.2,0) circle (1.7);
  \fill[orange!12] (3.2,0) circle (1.65);
  \fill[white] (3.2,0) circle (0.18);
  \draw[red!80!black, thick, ->] (4.35,0.65) arc[start angle=35,end angle=310,radius=1.32];
  \node at (3.2,0) {$\circ$};
  \node[below] at (3.2,-1.95) {小穿孔圆盘 $0<|q_h|<\varepsilon$};
  \node[right] at (3.35,0.05) {$q_h=0$};

  % filling point
  \draw[->, very thick] (5.2,0) -- (6.6,0);
  \node[above] at (5.9,0.15) {补上穿孔};
  \draw[thick, rounded corners=10pt] (7.0,-1.35) rectangle (9.0,1.35);
  \fill[red!80!black] (8.0,0) circle (2.4pt);
  \node[below] at (8.0,-1.65) {加入尖点后的坐标盘};
  \node[above, align=center] at (8.0,1.8) {把 $q_h=0$ 也纳入后\\尖点就像普通复点};
\end{tikzpicture}
\caption{尖点附近的局部模型。宽度为 $h$ 的条带在 $q_h=e^{2\pi i\tau/h}$ 下变成小穿孔圆盘；把穿孔点 $q_h=0$ 补上，就得到紧化后的尖点。}
\label{fig:cusp-q-coordinate}
\end{figure}


\begin{theorem}
\label{thm:x-gamma-compact-riemann}
把每个尖点 $s$ 对应的穿孔圆盘用坐标 $q_s$ 填上中心点 $q_s=0$，
便得到一个紧黎曼面
\[
  X(\Gamma)=\Gamma\backslash\HH^*.
\]
换言之，$X(\Gamma)$ 是 $Y(\Gamma)$ 的紧化，而每个新加进去的点正对应一个尖点。
\end{theorem}

\begin{proof}[思路]
先看尖点 $\infty$。
取足够大的 $T$，考虑条带
\[
  S_{h,T}=\{\tau\in\HH:\ |\Re\tau|\le h/2,\ \Im\tau>T\}.
\]
由于 $\tau\mapsto\tau+h$ 属于稳定子，$S_{h,T}$ 在商空间里描述了尖点附近全部点。
映射 $q_h=e^{2\pi i\tau/h}$ 把它送到穿孔圆盘
$0<|q_h|<e^{-2\pi T/h}$，并且是双全纯的。
于是只要把 $q_h=0$ 补进去，就得到一个真正的坐标盘。

对一般尖点 $s$，先用某个 $\sigma\in\SL_2(\Z)$ 把它搬到 $\infty$，
再做同样的构造即可。
因为有限指数子群只有有限多个尖点，
所以我们只补有限多个点；
而去掉这些尖点邻域后，剩下部分来自截断基本域的商，是紧的。
故整体得到一个紧黎曼面。
\end{proof}

\begin{proposition}[$\Gamma_0(N)$ 的尖点个数]
\label{prop:number-of-cusps-gamma0N}
$\Gamma_0(N)$ 的尖点个数为
\[
  \nu_\infty(N)=\sum_{d\mid N}\varphi\bigl(\gcd(d,N/d)\bigr).
\]
特别地，$\SL_2(\Z)$ 只有一个尖点，即 $\infty$。
\end{proposition}

\begin{proof}[说明]
在\cref{sec:congruence-subgroups} 里我们已经见过：
尖点就是 $\Gamma$ 在 $\mathbb{P}^1(\Q)$ 上的轨道。
对 $\Gamma_0(N)$，这些轨道可由分母数据分类，
最终恰好得到上式。
我们这里把它当作一个已知而非常有用的计数公式；
在亏格计算中，$\nu_\infty(N)$ 会直接进入 Riemann--Hurwitz 公式。
\end{proof}

\begin{example}[$X(1)$、$X_0(2)$ 与 $X_0(3)$ 的尖点]
\label{ex:cusps-small-levels}
\begin{enumerate}
  \item 对 $\SL_2(\Z)$，矩阵
    $S=\begin{pmatrix}0&-1\\1&0\end{pmatrix}$ 把 $\infty$ 送到 $0$，
    因而所有有理边界点都与 $\infty$ 等价，所以只有一个尖点。
  \item 对 $\Gamma_0(2)$，$\infty$ 与 $0$ 不再等价，故尖点恰有
    \[
      \{\infty,0\}.
    \]
    计数公式也给出
    $\varphi(1)+\varphi(1)=2$。
  \item 对 $\Gamma_0(3)$ 同样有两个尖点：
    \[
      \{\infty,0\}.
    \]
    因为 $3$ 是素数，公式仍给出 $2$。
\end{enumerate}
\end{example}

\begin{remark}
把尖点补进去以后，模形式在尖点处的 $q$-展开就有了几何含义：
"在尖点处全纯"，就是在新加进来的那个复点处全纯；
"是 cusp form"，就是在那个点处取值为 $0$。
下一节计算亏格时，尖点不仅是补上的点，
还是覆盖映射分歧的来源之一。
\end{remark}

\textbf{尖点的算术解释。}
在复几何里，尖点看起来只是补上的边界点；
但在算术几何里，它们对应的是\textbf{退化椭圆曲线}。
更准确地说，尖点附近的参数 $q_s$ 不只是复坐标，
也是 Tate 曲线族的模参数。

\begin{proposition}[Tate 曲线给出尖点的局部模型]
\label{prop:tate-curve-cusp-model}
设 $0<|q|<1$，并记 $q=e^{2\pi i\tau}$。
则 Weierstrass 方程
\[
  E_q:\ y^2+xy=x^3-\frac{E_4(q)}{48}x-\frac{E_6(q)}{864}
\]
定义一条椭圆曲线；在复解析上，它满足
\[
  E_q(\C)\cong \C^\times/q^\Z.
\]
当 $q\to 0$ 时，$E_q$ 退化为一个节点三次曲线；
从乘法型统一化的角度看，这正是把 $\mathbb{G}_m/q^\Z$ 的圆环极限捏成一个节点。
\end{proposition}

\begin{proof}[说明]
Tate 的 $q$-展开理论给出判别式
\[
  \Delta(q)=q\prod_{n\ge 1}(1-q^n)^{24},
\]
所以在 $0<|q|<1$ 上它从不为零，因而 $E_q$ 的 generic fiber 的确是光滑椭圆曲线。
另一方面，解析统一化 $\C^\times/q^\Z$ 说明：
$E_q$ 本质上是一个乘法群按 $q$ 的幂次缩放取商得到的复环面。
随着 $q$ 趋于 $0$，这个圆环被拉得越来越长；
在紧化极限中，唯一保留下来的奇性就是一个节点。
\end{proof}

\begin{figure}[ht]
\centering
\begin{tikzpicture}[font=\small, >=Stealth]
  \begin{scope}[xshift=-3.8cm]
    \draw[thick, blue!70!black] plot[smooth cycle] coordinates {(-1.45,0.0) (-0.75,0.95) (0.45,1.0) (1.35,0.2) (0.75,-0.95) (-0.55,-1.0)};
    \draw[thick, blue!45!black] plot[smooth cycle] coordinates {(-0.55,0.0) (-0.28,0.36) (0.18,0.38) (0.52,0.08) (0.28,-0.34) (-0.2,-0.38)};
    \node at (0,1.55) {$0<|q|<1$};
    \node at (0,-1.55) {$E_q(\C)\simeq \C^\times/q^\Z$};
  \end{scope}

  \draw[->, very thick] (-1.1,0) -- (1.1,0);
  \node[above] at (0,0.15) {$q\to 0$};

  \begin{scope}[xshift=3.8cm]
    \draw[thick, orange!80!black] (-1.45,-0.9) .. controls (-0.55,0.55) and (0.55,0.55) .. (1.45,-0.9);
    \draw[thick, orange!80!black] (-1.45,0.9) .. controls (-0.55,-0.55) and (0.55,-0.55) .. (1.45,0.9);
    \filldraw[orange!90!black] (0,0) circle (1.5pt);
    \node at (0,1.55) {节点纤维};
    \node at (0,-1.55) {尖点处的退化极限};
  \end{scope}
\end{tikzpicture}
\caption{Tate 曲线的退化示意。对 $0<|q|<1$，$E_q$ 由乘法型统一化 $\C^\times/q^\Z$ 给出；当 $q\to 0$ 时，圆环被拉长并最终捏成一个节点，得到尖点处的局部模型。}
\label{fig:ch04-tate-degeneration}
\end{figure}


\textbf{这与 $X_0(N)$ 的尖点有什么关系？}
若 $s$ 是 $X_0(N)$ 的一个尖点，宽度为 $h$，则局部参数
\[
  q_s=e^{2\pi i\sigma^{-1}\tau/h}
\]
正是相应 Tate 曲线的参数。
于是尖点附近的点不再只是“$\Im\tau$ 很大”的解析极限，
而是参数化带有一个 $N$-级循环子群的 Tate 曲线 $(E_{q_s},C)$；
当 $q_s=0$ 时，这个族退化为广义椭圆曲线的节点纤维。
CSS 在讨论整模型、坏约化和 $p$-进分析时反复利用的正是这一点：
尖点是模空间里真正的算术边界，
而 Tate 曲线给出了这个边界最自然的局部模型。
